\chapter{Lancer, tester, déboguer, surveiller}

Construire n'est qu'une des vingt-sept commandes. Ce chapitre couvre les quatre
qui servent tous les jours après.

\section{\texttt{run} --- lancer sans connaître le chemin}

\begin{nkcode}{}
jenga run --config Release
jenga run --project MonJeu --config Debug
jenga run -- --mode=solo --verbose
\end{nkcode}

\begin{nkretenir}
\textbf{Le double tiret sépare vos arguments de ceux de Jenga.} Tout ce qui suit
\texttt{-{}-} est passé au programme.

Sans lui, \texttt{jenga run -{}-verbose} rend Jenga bavard et ne transmet rien.
C'est une convention courante, et elle se rappelle par ce symptôme : « mon
programme ignore ses options ».
\end{nkretenir}

\begin{nkcode}{Capture reelle}
  EXECUTION  --  Bonjour.exe
     ...\Build\Bin\Release-Windows\Bonjour\Bonjour.exe

Application console sans terminal : ouverture d'une console dediee.
\end{nkcode}

\begin{nkretenir}
\textbf{Cette dernière ligne explique une confusion fréquente.} Une application
de type console lancée sans terminal attaché obtient une console à elle --- et sa
sortie n'apparaît donc pas dans \emph{votre} fenêtre.

Si vous capturez la sortie d'un programme par un script, vous ne verrez rien.
\emph{Ce n'est pas que le programme est muet ; c'est qu'il parle ailleurs.}
\end{nkretenir}

\section{\texttt{test} --- et ce qu'il répond quand il n'y a rien}

\begin{nkcode}{jenga test -- capture reelle, code de sortie 1}
No test projects found.
\end{nkcode}

\begin{nkretenir}
\textbf{Code de sortie 1, message explicite.} C'est le bon comportement : ne rien
trouver n'est pas un succès. Une intégration continue qui lance \texttt{jenga
test} sur un dépôt sans tests doit rougir, pas passer.

\emph{« Zéro test exécuté » et « zéro test en échec » se ressemblent au premier
coup d'\oe il et n'ont rien à voir.}
\end{nkretenir}

\begin{nkcode}{Declarer une suite de tests}
with project("MesTests"):
    testsuite()
    language("C++")
    files(["tests/**.cpp"])
    dependson(["MaBibliotheque"])
\end{nkcode}

\begin{nkretenir}[Le cadre de test est fourni par Jenga]
Jenga embarque \textbf{Unitest} --- \nkfile{Jenga/Unitest/} --- et le workspace
peut l'activer :

\begin{nkcode}{}
with unitest() as u:
    u.Precompiled()
\end{nkcode}

\alerte{} \textbf{C'est une confusion mesurée sur ce dépôt} : une recherche du cadre de
test dans les sources de Nkentseu avait conclu « ce cadre n'existe pas ». Il
existe --- \emph{dans Jenga}. La recherche était exhaustive et sa racine était
fausse.

\emph{Une recherche exhaustive ne vaut que ce que vaut sa racine.}
\end{nkretenir}

\begin{nkretenir}
Deux réglages coupent les tests, et il faut savoir qu'ils existent avant de
conclure qu'une suite « ne tourne pas » :

\begin{center}
\begin{tabular}{@{}ll@{}}
\toprule
\texttt{disableunittestcompilation()} / \texttt{dutc()} & ne les compile plus \\
\texttt{disableunittestexecution()} / \texttt{dute()} & les compile, ne les lance pas \\
\bottomrule
\end{tabular}
\end{center}

\textbf{Le second est le plus trompeur} : tout se construit, rien ne s'exécute.
Depuis la version 2.4.0, \texttt{test} et \texttt{run} acceptent
\texttt{-{}-force} pour lever ces deux verrous.
\end{nkretenir}

\begin{nkpiege}
\textbf{Un enchaînement \texttt{A || B} attribue la sortie de B à A.} Mesuré :
un développeur lance \texttt{jenga test -{}-project X || jenga build -{}-target X},
lit \texttt{6/6 --- SUCCESS}, et conclut que \texttt{jenga test} annonce un succès
sur une suite qu'il n'exécute pas.

Relancées séparément, \texttt{jenga test} répondait \textbf{en rouge} et le
succès venait du \emph{build}, qui avait raison de réussir. Un défaut d'outil
allait être inscrit dans un document partagé.

\textbf{En diagnostic : une commande, une lecture.}
\end{nkpiege}

\section{\texttt{gdb} --- déboguer}

\begin{nkcode}{}
jenga gdb --config Debug
jenga debug --config Debug     # meme commande
\end{nkcode}

\begin{nkretenir}
Jenga lance le débogueur sur le bon exécutable, avec la bonne configuration.
\textbf{Construisez en Debug} : en Release, les symboles sont absents
(\texttt{symbols(False)}) et vous déboguerez à l'aveugle.

Pour un plantage non interactif, la forme qui donne le plus en une fois :

\begin{nkcode}{}
gdb --batch -ex "run <args>" -ex "bt full" --args <exe>
\end{nkcode}
\end{nkretenir}

\section{\texttt{watch} --- reconstruire à chaque enregistrement}

\begin{nkcode}{}
jenga watch
jenga w
\end{nkcode}

\begin{nkretenir}
\texttt{watch} surveille vos sources et reconstruit à chaque modification. Il
exige \texttt{watchdog} --- la seule dépendance Python vraiment optionnelle de
Jenga.

\textbf{Il hérite de la décision par dates du chapitre 9}, et cette fois à votre
avantage : un enregistrement met la date à \emph{maintenant}, donc la
recompilation part.
\end{nkretenir}

\begin{nkpiege}
\textbf{Ne laissez pas \texttt{watch} tourner pendant une mesure.} Il reconstruit
au moindre changement --- y compris ceux d'un autre outil --- et vous mesureriez
un binaire différent de celui que vous croyez.

Même règle que pour toute campagne : on ne répare pas la tuyauterie pendant que
l'eau coule.
\end{nkpiege}

\section{\texttt{clean} et \texttt{rebuild}}

\begin{nkcode}{jenga clean -- capture reelle, extrait}
Removed ...\Build\Obj\Release-Windows\Bonjour\src_main.obj
Removed ...\Build\Bin\Release-Windows\Bonjour\Bonjour.exe
\end{nkcode}

\begin{nkretenir}
\textbf{\texttt{clean} énumère ce qu'il supprime}, ce qui est plus utile qu'il
n'y paraît : la liste montre \emph{quelle} configuration et \emph{quelle}
plateforme sont nettoyées.

\texttt{jenga clean} sans argument ne nettoie pas nécessairement tout : les
autres configurations restent. C'est voulu --- et c'est la raison pour laquelle
un « j'ai pourtant nettoyé » ne clôt pas une enquête.
\end{nkretenir}

\begin{nkretenir}
Notez aussi le nom de l'objet : \texttt{src\_main.obj}. Le chemin de la source
est aplati dans le nom --- \texttt{src/main.cpp} devient \texttt{src\_main}.

C'est ce qui permet à deux fichiers du même nom dans deux dossiers de coexister,
et c'est utile à connaître quand on cherche si \emph{son} fichier a bien été
compilé (chapitre 9).
\end{nkretenir}

\begin{nkbilan}
\texttt{run} retrouve l'exécutable et passe vos arguments après \texttt{-{}-} ;
une application console sans terminal parle dans une console à elle, ce qui
explique bien des « programmes muets ». \texttt{test} rend un code non nul quand
il ne trouve rien --- « zéro exécuté » n'est pas « zéro en échec » --- et le cadre
Unitest est fourni par Jenga, pas par votre dépôt. \texttt{gdb} veut du Debug,
\texttt{watch} veut \texttt{watchdog} et ne doit pas tourner pendant une mesure.
Enfin \texttt{clean} énumère, et ne nettoie que la configuration visée.
\end{nkbilan}

\begin{nkquestions}
\begin{enumerate}
\item À quoi sert \texttt{-{}-} dans \texttt{jenga run} ?
\item Que rend \texttt{jenga test} sur un projet sans tests, et pourquoi est-ce
      le bon comportement ?
\item Quelle différence entre \texttt{dutc} et \texttt{dute} ?
\item Pourquoi ne faut-il pas laisser \texttt{watch} tourner pendant une mesure ?
\item D'où vient le nom \texttt{src\_main.obj} ?
\end{enumerate}
\end{nkquestions}

\begin{nkpratique}
Sur votre projet :

\begin{enumerate}
\item ajoutez une suite de tests avec un cas qui passe et un qui échoue ;
\item lancez \texttt{jenga test} et relevez le \textbf{code de sortie} ;
\item corrigez le cas fautif, relancez, relevez à nouveau ;
\item posez \texttt{dute()} dans le workspace et relancez --- que se passe-t-il,
      et le code de sortie change-t-il ?
\end{enumerate}

Le quatrième point est le plus important : c'est celui qui rend un test
silencieux.
\end{nkpratique}

\begin{nkdemo}
Prouvez qu'un enchaînement masque l'origine d'une sortie.

Lancez \texttt{jenga test || jenga build} sur un projet \emph{sans} tests, et
relevez ce que vous lisez à l'écran.

Puis lancez les deux commandes \textbf{séparément} et relevez chaque code de
sortie.

\textbf{Ce que la démonstration établit} : la première forme affiche un succès,
et rien n'indique qu'il vient de la seconde commande. Vous auriez pu conclure que
\texttt{test} réussit.

C'est le mécanisme exact qui a fait accuser Jenga à tort sur ce dépôt --- et le
défaut allait être inscrit dans un document partagé.
\end{nkdemo}

\begin{nklogique}
Une suite de tests « ne s'exécute plus » depuis quelques jours, sans que personne
n'ait touché aux tests.

\begin{enumerate}
\item Énumérez quatre causes, en distinguant celles qui empêchent la
      \emph{compilation} de celles qui empêchent l'\emph{exécution}.
\item Pour chacune, dites quel signe la trahit --- et laquelle ne laisse
      \emph{aucun} signe dans la sortie de \texttt{build}.
\item Quel contrôle, ajouté à votre intégration continue, rendrait cette dernière
      impossible à ignorer ?
\end{enumerate}
\end{nklogique}
